\section{How the Church Learned to Wait}\label{sec:history}

When Western Christians began keeping the Nativity on 25~December in the
fourth century, no universal preparation preceded the feast. Over four
centuries, churches in Spain, Gaul, Milan, and Rome each learned to wait in
their own way---through twenty-one days of unbroken church attendance,
through a six-week fast after St.~Martin's Day, and in Rome through a
December fast of harvest thanksgiving and almsgiving that gradually grew
into a four-Sunday season.

\subsection{Spain, Gaul, and the Non-Roman Western Rites}

The earliest Western council to fix a winter watch met at Saragossa in
Spain (dated 380, or 379 in some chronologies). Twelve bishops gathered
under Pope St.~Damasus against the Priscillianists, whose followers were
withdrawing from parish churches during the winter feasts. Canon~4 ordered
that for twenty-one consecutive days, from 17~December (\textit{a 16.\
Kalendas Januarii}) to Epiphany on 6~January (\textit{8.\ Idus Januarii}),
no one might absent himself from church, hide at home, retire to country
estates, flee to the mountains, or walk barefoot; all had to assemble at
church (\textit{concurrere ad ecclesiam}). Saragossa mandated attendance
rather than fasting and never named the interval ``Advent,'' yet its
opening date, 17~December, later marked the opening of the Roman Great
O~Antiphons.

North of the Pyrenees, Christmas preparation took a penitential form
anchored to St.~Martin's Day (11~November). In the closing chapter of his
\textit{History of the Franks} (X.31), written in the late sixth century,
St.~Gregory of Tours preserved the yearly fasts ordered by his
late-fifth-century predecessor, Bishop Perpetuus of Tours, including the
clause: ``From the burial of our lord Martin to the Nativity of our Lord,
three fasts each week'' (\textit{De depositione domni Martini usque natale
Domini terna in septimana ieiunia}). One early manuscript omits that
clause while another preserves it (Appendix~\ref{app:about}).

Two sixth-century Gallican councils confirm how firmly the December fast
took root. The Second Council of Tours (567) bound monks specifically
(\textit{a monachis}) to fast three days a week from September through
November and daily throughout December until Christmas (\textit{de Decembre
usque natale Domini omni die ieiunent}). Soon afterward, the First Council
of M\^acon (dated between 581 and 583) met at King Guntram's summons with
twenty-one bishops attending. Canon~9 extended a Monday, Wednesday, and
Friday fast from St.~Martin's Day to Christmas across the church, ordered
Mass celebrated in the Lenten manner (\textit{sacrificia quadragensimali
debeant ordine caelebrari}), and directed that the canons be read aloud so
none could plead ignorance. Out of this six-week Gallican discipline grew
the name ``St.~Martin's Lent.''

Other Western rites shaped their own cycles. At Milan, the Ambrosian
liturgy opened after St.~Martin's Day and kept six Sundays of Advent,
preserved in the eleventh-century \textit{Manuale Ambrosianum} as
\textit{Dominica I} through \textit{Dominica VI de Adventu}. In ancient
Spain, the earliest Mozarabic manuscripts count no more than five Sundays
of Advent, while the later printed Mozarabic Missal and Breviary count six.

\subsection{Rome: From the Tenth-Month Fast to Four Sundays}

Rome followed a distinct course. Between 440 and 461, Pope St.~Leo the
Great preached every December on the \textit{jejunium decimi mensis}, the
Fast of the Tenth Month (December in the old Roman calendar starting in
March). In Sermons 12, 13, 16, 17, and~19, Leo joined this winter fast to
Lent, Pentecost, and September as a fourfold yearly rhythm (\textit{Sermon
19.2}):

\begin{quote}
Accordingly we keep the spring fast in Lent, the summer fast at
Whitsuntide, the autumn fast in the seventh month, and the winter fast in
this which is the tenth month, knowing that there is nothing unconnected
with the Divine commands, and that all the elements serve the Word of God
to our instruction, so that from the very hinges on which the world turns,
as if by four gospels we learn unceasingly what to preach and what to do.
\end{quote}

For Leo, the December fast offered thanksgiving after ``the ingathering of
the crops was complete'' so each believer might be ``more abstinent in
himself and more open-handed towards the poor'' (\textit{Sermon 16.2}). He
summoned Rome to fast on Wednesday and Friday, to turn fasting into food
for the poor (\textit{Fiat refectio pauperis, abstinentia jejunantis},
\textit{Sermon 13}), and to keep Saturday vigil at St.~Peter's. Leo never
named a liturgical season ``Advent''; even when \textit{Sermon 19.1} cites
Luke 21:34 on preparing for the \textit{adventus} of God's kingdom, he
speaks of the coming Day of the Lord. Only later, after Leo's sermons
entered Advent Matins, did modern translators occasionally insert
``Advent'' where his Latin named only the tenth month.

Early Roman sacramentaries trace the tightening from five Masses to four
Sundays. The Old Gelasian Sacramentary---an eighth-century Frankish copy
preserving seventh-century Roman prayers---holds five Advent Masses under
\textit{Orationes de adventum Domini} (Book~II, LXXX--LXXXIV). It does not
number them as Sundays, though companion eighth-century books from Rheinau
and St.~Gall assign them to the five Sundays before Christmas.

The \textit{Homilies on the Gospels} of Pope St.~Gregory the Great
(preached around 590--592) catch the Roman arrangement mid-course.
Homily~1, on the signs in sun and moon (Luke 21:25--33), bears the heading
\textit{Dominica secunda Adventus Domini}---the Second Sunday of
Advent---at St.~Peter's (a Lateran copy places it at St.~Mary Major
\textit{ad praesepe}); only later did that Gospel shift to the First
Sunday. Homily~20 (on Luke 3:1--6) is headed \textit{Sabbato quatuor
temporum ante Natalem Christi}, for Ember Saturday before Christmas, though
four manuscripts omit that Ember rubric.

By the late eighth century, the papal sacramentary sent by Pope Hadrian~I
to Charlemagne---the Gregorian \textit{Hadrianum}---settled on four
Sundays: \textit{Dominica Prima}, \textit{Dominica Secunda}, and
\textit{Dominica III ad Sanctum Petrum}, followed by Ember Wednesday,
Friday, and Saturday at St.~Peter's, and \textit{Die Dominica vacat}, a
``vacant'' fourth Sunday. No daytime stational Mass stood there because the
Roman church spent Saturday night in the twelve-lesson Ember vigil and
ordinations until Sunday dawn. Beneath that vacant heading, however, the
\textit{Hadrianum} already placed the prayer opening \textit{Excita domine
potentiam tuam et ueni} (from the Old Gelasian), which became the Fourth
Sunday Collect \textit{Excita, quaesumus, Domine, potentiam tuam, et veni:
et magna nobis virtute succurre} once a daytime Mass was assembled from
existing Advent chants and readings.

\subsection{Medieval Liturgists and the Sound of Advent}

When Roman books spread across Charlemagne's empire, Frankish clergy found
both counts on their lecterns. Writing around 823--835, Amalarius of Metz
noted (\textit{Liber officialis} III.40) that older missals and
lectionaries began five weeks before Christmas (\textit{Hebdomada quinta
ante Natalem Domini}), while the ``Gregorian'' Missal (\textit{Missalis qui
vocatur Gregorialis}) and Antiphoner appointed four weeks (\textit{per
quatuor hebdomadas}). Amalarius also recalled (\textit{Vidi tempore
prisco}) that in Advent \textit{Gloria in excelsis Deo} was omitted---and
in some places deacons laid aside dalmatics---to signify Old Testament
waiting.

Twelfth- and thirteenth-century liturgists---Rupert of Deutz (\textit{De
divinis officiis} III.1--2), Honorius Augustodunensis (\textit{Gemma
animae} III.1), Sicard of Cremona (\textit{Mitrale} V.1), John Beleth, and
William Durandus---kept explaining both the five-week books and the
character of the four-week Roman season. They also took up a question every
observant worshipper asks: if \textit{Gloria in excelsis} and \textit{Te
Deum} fall silent in Advent, why does \textit{Alleluia} stay?

In his \textit{Summa on the Ecclesiastical Offices} (ch.\ LXIV), John
Beleth wrote that \textit{Te Deum} is withheld so it may be resumed with
greater ardor on the feast, or ``because He whom we await is not yet
present'' (\textit{quia nondum pr\ae sens est quem exspectamus}). Along
with \textit{Gloria in excelsis} and (in medieval usage) \textit{Ite, missa
est}, it falls quiet; ``yet because this season is partly of rejoicing''
(\textit{cum tempus istud partim sit l\ae titi\ae}), Beleth observed,
``the \textit{Alleluia} is not silenced as in Septuagesima, which is a
season wholly of mourning'' (\textit{totum m\oe roris}). A century later,
William Durandus (\textit{Rationale} VI.2, c.\ 1286) echoed him:
\textit{Gloria in excelsis} belongs to Christmas Night and is saved to be
sung more eagerly when that night arrives; as for why \textit{Alleluia}
stays amid fasting and suspended weddings, this fast is ``partly of
rejoicing, and partly of mourning.'' (Durandus also repeated a medieval
legend attributing three full weeks and an incomplete fourth week of Advent
to St.~Peter---medieval tradition rather than first-century record.)

\subsection{From Trent to the 1962 Horizon}

The Council of Trent and St.~Pius~V fixed the Roman order for the Latin
Church. On 11~November 1563, Session~XXIV (\textit{De reformatione
matrimonii}, ch.\ 10) renewed the closed time forbidding solemn weddings
``from the Advent of our Lord Jesus Christ until the day of the Epiphany''
(later shortened to Christmas Day in the 1917 Code of Canon Law; see
Section~\ref{sec:penance-law}). On 14~July 1570, St.~Pius~V promulgated
the restored \textit{Missale Romanum} by \textit{Quo primum}, establishing
four-Sunday Roman Advent across the West while sparing local rites over two
hundred years old, preserving Milan's six-Sunday Advent and Mozarabic
Spain.

Over the centuries, lay Advent fasting gradually relaxed outside the Ember
Days and Christmas Vigil---leaving nineteenth-century remnants such as
Wednesday and Friday fasting in England and Friday fasting in most American
dioceses in the \textit{Manual of Prayers} ordered by the Third Plenary
Council of Baltimore (1884; published 1889)---while the liturgical prayers
stayed steady. Three acts between 1958 and 1962 gave the 1962 Roman Rite
its exact shape:

\begin{itemize}
\item \textbf{\textit{De musica sacra} (3~September 1958).} The Sacred
Congregation of Rites (nn.\ 81~a, 83) barred solo organ and instruments
from First Vespers of Advent~I through None of the Christmas Vigil (outside
Eucharistic Benediction), while allowing instruments on holy days and major
feasts (except Sundays), solo organ or harmonium on Gaudete Sunday, and
organ or harmonium at Mass and Vespers \textit{solummodo ad cantum
sustentandum}, solely to sustain singing.

\item \textbf{\textit{Rubricarum instructum} and \textit{Codex Rubricarum}
(25--26~July 1960; effective 1~January 1961).} Approved by Pope St.~John~XXIII
and promulgated the next day, the Code fixed Advent from First Vespers of
Advent~I through None inclusive of the Vigil (RG n.\ 71), with Advent~I on
or nearest to 30~November (RG n.\ 20). Keeping the four Sundays as
First-Class Sundays (RG n.\ 11), it ranked Ember Days and 17--23~December
as Second-Class ferias and earlier weekdays as Third-Class (RG nn.\ 24--25),
and ruled that 24~December on Sunday replaces Advent~IV without
commemoration (RG n.\ 30~a). It restored \textit{Ite, missa est} at Advent
Mass unless a procession follows (RGMR n.\ 507), abolished folded chasubles
(\textit{planetae plicatae}) and the broad stole (\textit{stola latior}) so
ministers wear dalmatic and tunicle (RG n.\ 137), restricted rose vestments
in Advent to Gaudete Sunday's own Office and Mass alone (RG n.\ 131),
incorporated the 1955 suppression of standing seasonal suffrages and extra
collects, and simplified Sunday Matins to one nocturn of nine psalms and
three lessons---two from Isaiah and one from the Gospel homily (RGBR
nn.\ 162~a, 220).

\item \textbf{The 1962 \textit{Missale Romanum} (23~June 1962).} The Sacred
Congregation of Rites issued the typical edition harmonized with the 1960
Code (and on 13~November 1962 Pope St.~John~XXIII ordered St.~Joseph added
to the Roman Canon from 8~December 1962).
\end{itemize}

Ceremonial change never severed continuity in prayer. When a priest opens
Advent~I in the 1962 Missal with \textit{Excita, quaesumus, Domine,
potentiam tuam, et veni}, he prays the same petition standing at the head
of \textit{Dominica Prima} in the eighth-century \textit{Hadrianum}
(\textit{Excita domine quaesumus potentiam tuam et ueni}), inside a season
that still keeps the tenth-month Ember fast Pope St.~Leo preached in
fifth-century Rome.

\begin{longtable}{>{\raggedright\arraybackslash}p{0.18\linewidth}>{\raggedright\arraybackslash}p{0.26\linewidth}>{\raggedright\arraybackslash}p{0.48\linewidth}}
\toprule
\textbf{Date} & \textbf{Place and Source} & \textbf{Development in Advent} \\
\midrule
\endhead
380 (or 379) & Saragossa, Council, c.\ 4 & 21 days of church attendance (17~Dec--6~Jan) against Priscillianist withdrawal; no fast, no name ``Advent.'' \\
440--461 & Rome, St.~Leo, Sermons 12, 13, 16, 17, 19 & Tenth-month fast (\textit{jejunium decimi mensis}): Wed/Fri fast, alms, Sat vigil at St.~Peter's; never names ``Advent.'' \\
Late 5th c.\ (rec.\ 6th c.) & Tours, Bp.\ Perpetuus (\textit{Hist.\ Franc.} X.31) & Three weekly fasts from St.~Martin's Day to Christmas (clause omitted in one early manuscript). \\
567 & Tours, Second Council & Binds monks to fast daily in December until Christmas. \\
581--583 & M\^acon, First Council, c.\ 9 & King Guntram and 21 bishops decree Mon/Wed/Fri fast and Lenten Mass from Martinmas to Christmas. \\
c.\ 590--592 & Rome, St.~Gregory, \textit{Homilies} & Gospel homilies with Advent rubrics (Homily~1 headed for Sun~II; Homily~20 for Ember Sat). \\
8th c.\ MS (7th c.\ core) & Old Gelasian Sacramentary & Five Advent Masses (\textit{Orationes de adventum Domini}), assigned to five Sundays in companion copies. \\
Late 8th c. & Rome / Aachen, \textit{Hadrianum} & Four Sundays (\textit{Dominica I--III}, Ember week, \textit{Die Dominica vacat}) with \textit{Excita domine quaesumus}. \\
c.\ 823--835 & Metz, Amalarius, \textit{Lib.\ off.} III.40 & Contrasts five-week older books with four-week Gregorian Missal; notes omitted \textit{Gloria} and dalmatics. \\
11th--12th c. & Milan and Mozarabic Spain & \textit{Manuale Ambrosianum} has six Sundays; earliest Mozarabic MSS have five (six in print). \\
12th--13th c. & Beleth (\textit{Summa} LXIV), Durandus (\textit{Rat.} VI.2) & Explain silent \textit{Gloria} and \textit{Te Deum} with \textit{Alleluia}: partly rejoicing, partly mourning. \\
1563, 1570 & Trent (XXIV c.\ 10); St.~Pius~V & Renews closed time for solemn weddings (Advent--Epiphany); \textit{Quo primum} restores Roman Missal (200-year exemption). \\
3~Sept 1958 & Rome, SCR, \textit{De musica sacra} & Bars solo organ and instruments except major feasts, Gaudete Sunday, and solely to sustain singing. \\
25--26~July 1960 & Pope St.~John~XXIII, \textit{Codex Rubricarum} & Eff.\ 1~Jan 1961: bounds I~Vesp.\ I--Vigil None; \textit{Ite, missa est}; no folded chasubles; rose Gaudete only in Advent; 3-lesson Matins. \\
23~June 1962 & Rome, SCR, \textit{Missale Romanum} & Typical edition harmonized with 1960 Code (St.~Joseph added to the Canon on 8~Dec 1962). \\
\bottomrule
\end{longtable}
